\section{逐项任务与完整验收矩阵}
\subsection{实施任务 T00--T14}
本附录编号来自本次设计核查任务卡，与测试契约 T01--T16 不同。任务 T00--T09、T12、T13 已按默认静态 profile 完成；T10/T11 为条件能力；T14 仅完成前瞻协议。
\begin{longtable}{L{14mm}L{69mm}L{66mm}}
\caption{任务完成范围；不将条件项或协议冻结视为测量完成。}\\
\toprule 任务 & 内容 & 当前状态\\\midrule\endfirsthead
\caption[]{（续）}\\
\toprule 任务 & 内容 & 当前状态\\\midrule\endhead
T00 & 实际基线、旧修改/历史保护、特征化 & 默认静态范围完成。\\
T01 & 性质契约、状态/身份/分母与导出 schema & 默认静态范围完成。\\
T02 & 产品能力预检与功能检查注册 & 默认静态范围完成。\\
T03 & 独立产品服务和真实 CLI & 默认静态范围完成。\\
T04 & 结构化诊断、准备/等待输入和补项 & 默认静态范围完成。\\
T05 & 公共快照与分层身份 & 默认静态范围完成；私有反射非保证。\\
T06 & 全局证据归并与研究重新准入 & 已知投递事实范围完成。\\
T07 & 真实数据依赖、有限控制与图上限 & 默认静态范围完成，旧研究义务保留。\\
T08 & 守恒证书池、缓存及独立投递计费 & 默认静态范围完成；服务无必成证明预留。\\
T09 & 显式行范围视图与付费扩展 & 默认静态范围完成。\\
T10 & 私有 scratch 与生命周期 & 条件未启用。\\
T11 & 可信 native build & 条件未启用。\\
T12 & runtime 记账、只读恢复与取消/脱敏 & 启用静态范围及保留控制完成。\\
T13 & 完整 pytest、harness、真实 CLI 与矩阵 & 默认静态范围已验收。\\
T14 & 人工/模型增量协议 & 协议冻结；GP3/GP4 未测量。\\
\bottomrule\end{longtable}

\subsection{A01--A34 的验收}
\begin{longtable}{L{11mm}L{11mm}L{97mm}L{30mm}}
\caption{全部 34 个验收项。31 必需项通过，3 条件项保留未提供。每项测试路径和日志 SHA-256 在 acceptance.json/tasks.json。}\\
\toprule 项 & 任务 & 验证行为 & 结果\\\midrule\endfirsthead
\caption[]{（续）}\\
\toprule 项 & 任务 & 验证行为 & 结果\\\midrule\endhead
A01 & T00 & 实际基线、历史保护、完整套件与特征化 & 通过\\
A02 & T01 & 20 个运行/性质组合、非法状态和计数 & 通过\\
A03 & T01 & 精确 subject、性质、行、source/analysis 绑定 & 通过\\
A04 & T01 & 应用未评估和 remote/public false & 通过\\
A05 & T02 & 无 provider/witness 的静态能力 & 通过\\
A06 & T02 & 已注册自有检查缺依赖及显式恢复 & 通过\\
A07 & T03 & 实际 console 与 HTTP/socket/gate 拒绝 spy & 通过\\
A08 & T03 & 零候选与解析失败分母/产物 & 通过\\
A09 & T03 & 内部失败保留既有 proof 和局部产物 & 通过\\
A10 & T04 & F4 环境与 F1 准备区分，F1--F6 保留 & 通过\\
A11 & T04 & 等待输入、补项及依赖恢复，准备非完成 & 通过\\
A12 & T05 & 五项嵌套别名回归及源码正反控制 & 通过\\
A13 & T05 & 源码/helper/import/环境漂移与视图分离 & 通过\\
A14 & T08 & 别名与两个原 ID 区分，保留验证记账 & 通过\\
A15 & T06 & 晚到冲突两种顺序，slow/fast 重新准入 & 通过\\
A16 & T06 & 局部反事实不抹去完整独立支持 & 通过\\
A17 & T06 & 三个引用别名不增加强度 & 通过\\
A18 & T06 & UNKNOWN 不抹证明，helper 引用保留 & 通过\\
A19 & T07 & 不可用可选任务与独立静态任务，真实依赖阻塞 & 通过\\
A20 & T07 & 完成 UNRESOLVED 解锁，准备/error 不解锁 & 通过\\
A21 & T07 & 固定无效动作限额及无进展图终止 & 通过\\
A22 & T08 & 2500+1000+1000+1000 完整证书在 8192 内 & 通过\\
A23 & T08 & JSON 外壳/ID 计费，裁剪及局部错误不证明 & 通过\\
A24 & T08 & 缓存投递计 bytes，原执行唯一且保护预留 & 通过\\
A25 & T09 & 授权 helper 投递计费；失败扩展不准入 & 通过\\
A26 & T09 & 同路径不准入其他行范围，旧 E 预算保留 & 通过\\
A27 & T10 & 私有 scratch 能力未启用 & 未启用 / 未运行\\
A28 & T10 & 新 scratch 生命周期未启用，保留旧只读控制 & 未启用 / 未运行\\
A29 & T11 & 可信 native build 未启用 & 未启用 / 未运行\\
A30 & T12 & 每进程完整 gate、漂移及失败拒绝 & 通过\\
A31 & T07 & 八个实际任务，循环受图上限 16 约束 & 通过\\
A32 & T12 & 六个最后名额竞争者及畸形内容 usage 保留 & 通过\\
A33 & T12 & 半写 journal、未知 reservation、取消无后续调用 & 通过\\
A34 & T12 & 伪状态/源码 key 留数据私有，transport key 脱敏 & 通过\\
\bottomrule\end{longtable}
\source{audit/design_review_20260930/acceptance.json; audit/design_review_20260930/tasks.json}
\subsection{完整日志与源码摘要}
最终完整日志 SHA-256 为 \code{d3a2bc10a5e0d9480039c18486c74689d8a3970fa20a19ab4af445c7efd1c341}。报告收集的 230 个 Python 源文件聚合摘要为 \code{c1bd6bb77ff68a2449b6b7d8af33814e286c224fd733e960f7b0e5ac89b58fa8}；算法按排序路径加入路径 bytes、分隔零字节和每文件 SHA-256 bytes。该摘要用于报告工作树复核，与产品 source-bytes-v1 或研究 ValidationSubject 的摘要算法不同。

\subsection{保留的语义测试契约}
C1--C6 是源码语义的规范职责。下表来自当前测试契约；T01--T16 是语义回归编号，与上文实施任务卡不同。
\begin{longtable}{L{13mm}L{35mm}L{101mm}}
\caption{语义契约及其执行要求。}\\
\toprule 契约 & 主题 & 核心要求\\\midrule\endfirsthead
\caption[]{（续）}\\
\toprule 契约 & 主题 & 核心要求\\\midrule\endhead
C1 & 判断对象 & 候选、源码、位置、性质、范围及前提完整绑定；局部事实不升级为不匹配性质的证明。\\
C2 & 值与执行效果 & 区分已知值、来源、未知值，正常继续、终止、异常，及相关副作用与覆盖缺口。\\
C3 & 求值顺序与摘要 & 短路、条件表达式、调用实参与 helper 服从同一副作用规则；摘要需要身份和输入前提。\\
C4 & 路径合并 & 完整危险路径、全范围安全、未解析路径、同范围冲突分别处理；重复条件关联不能制造假路径。\\
C5 & 证据准入 & 全部相关观察约束判决；调用顺序、引用角色和静态扫描增删不能替代语义事实。\\
C6 & 收尾与依赖 & 终止不降低证据门槛；任务结束不自动等于验证完成；非法协议不伪装成普通弃权。\\
\bottomrule
\end{longtable}
\begin{longtable}{L{13mm}L{46mm}L{90mm}}
\caption{测试契约要求的验收类别。当前主链核验不表示所有组合已穷尽。\citep{contract}}\label{tab:acceptance}\\
\toprule 类别 & 检查主题 & 关键行为与正控制\\\midrule\endfirsthead
\caption[]{（续）}\\
\toprule 类别 & 检查主题 & 关键行为与正控制\\\midrule\endhead
T01 & caller/helper 范围 & 不借 helper 局部安全抵消 caller 危险；双方受支持安全时仍保留安全判断。\\
T02 & 位置与切片 & 不同 sink、偏移、装饰器及 helper 不错借身份；准确位置正例可通过。\\
T03 & 表达式副作用 & 条件、短路、比较、海象和参数位置的相关操作被保留或明确未解析。\\
T04 & 纯度与死分支 & 确定死分支和未调用延迟体不虚构执行；未知对象协议不默认纯净。\\
T05 & helper 执行效果 & 参数、别名、全局和函数重绑定缺口不被未使用返回值抹去；独立完整危险事实保留。\\
T06 & 解包与终止 & 全成功、全失败、成功与失败/未知组合、嵌套及异常边界保持正确后继。\\
T07 & 输入来源 & 未知值不等于已证明输入；错来源不套用摘要；具体前提保持正例。\\
T08 & 消费上下文 & 正常保护、后续拼接、二次解释与非 shell 消费区别处理。\\
T09 & 条件关联 & 相同条件重复检查不生成不可能路径；绑定改变和未声明平台保持缺口。\\
T10 & 权限覆盖 & 不同目标和时间顺序不抵消；位模式证明不升级为竞态证明。\\
T11 & 引用与取证不变性 & 角色、省略、静态扫描和独立读取顺序不能改变同一语义事实的约束。\\
T12 & 路径和证明冲突 & 不同独立路径与同范围互斥证明分别处理；真实危险路径正例仍可用。\\
T13 & 身份与投递 & 完整且匹配的证明可用；错 subject、性质、摘要、引用及裁剪/抑制/错误被拒。\\
T14 & 收尾和用量 & 普通与末步阶段同一证据门槛，零尝试不伪造成票，全部实际请求记账。\\
T15 & 验证完成与依赖 & 未完成必需验证器阻塞直接和传递依赖；已完成未决与未完成分别记录。\\
T16 & 多任务与 quorum & 同一专家一票，弱票不能凑证明，固定票早停与全量标签一致。\\
\bottomrule
\end{longtable}
